\chapter{Introduction}
\label{ch:introduction}

This monograph is organized around a single question: when does structure
survive transformation, and when does it not? The question has appeared
throughout this research program under several names --- reconstruction,
repair, restorability, continuation --- and has been approached from
several directions, including a formal treatment of admissibility and
irreversibility, a foundational account of observability as restorability,
and an extended study of continuation across unequal timescales. What
follows is not a new direction but a consolidation: four theorems, each
of which has appeared informally or partially elsewhere in this program,
stated here as precise results and derived directly from the admissibility
apparatus rather than argued for case by case.

The apparatus itself is simple to state. An admissibility structure $R$
determines, for any state $x$, a reachable set $\reach{x}$ of states $y$
such that $x \adm{R} y$. Nothing about $R$ is assumed beyond this: it is
a relation, not a metric, and it need not be symmetric, transitive, or
total. Historical work in this program has shown that a great deal of
structure can be recovered from this minimal starting point once it is
allowed to vary with time, $R_t$, and once transitions between admissible
structures are themselves treated as admissible or inadmissible events
rather than as external clock ticks.

The four theorems developed in this monograph correspond to four
questions one can ask of this apparatus. First: when are two
representations, possibly built from entirely different physical or
computational substrates, the same object as far as $R$ is concerned?
This is Representation Invariance, the subject of Chapter~\ref{ch:invariance}.
Second: what does it mean for a trajectory through admissibility structure
to be coherent, and is coherence an additional postulate or a consequence
of the relation itself? This is Admissibility Preservation, the subject
of Chapter~\ref{ch:preservation}. Third: under what conditions, and in
what precise sense, does a scalar measure of inadmissibility decrease
along a trajectory, and how does this relate to --- and how must it be
distinguished from --- the ordering relation of monotonic continuation
developed elsewhere in this program? This is Monotonic Relaxation, the
subject of Chapter~\ref{ch:relaxation}. Fourth: when is it structurally
impossible, not merely difficult, to recover which of two histories
actually occurred from the evidence either would have left behind? This
is Non-Identifiability, the subject of Chapter~\ref{ch:nonidentifiability}.

A methodological note is appropriate before proceeding. Each of the four
theorems could be, and in earlier work occasionally has been, treated as
an independent claim requiring its own justification. The approach taken
here is different. Every theorem below is derived from definitions
already established elsewhere in this program --- the admissibility
relation, the reachable set, the closure functional --- so that what is
new in this monograph is not the vocabulary but the demonstration that
four apparently distinct phenomena are, in fact, four views of one
relation. Where a theorem's content is close to definitional, this is
stated plainly rather than disguised as a deeper result; where a theorem
requires a genuine argument, as with Admissibility Preservation and
Monotonic Relaxation, the argument is given in full and its conditions
of validity are stated explicitly rather than left implicit.
